\documentclass[11pt,a4paper]{article}

\usepackage[T1]{fontenc}
\usepackage[utf8]{inputenc}
\usepackage{amsmath,amssymb}
\usepackage{geometry}
\usepackage{hyperref}
\usepackage{xcolor}
\usepackage{booktabs}
\usepackage{parskip}

\geometry{a4paper, left=3cm, right=3cm, top=3cm, bottom=3cm}

\hypersetup{colorlinks=true, linkcolor=blue!60!black, urlcolor=blue!60!black}

\definecolor{notecolor}{rgb}{0.25, 0.35, 0.60}
\definecolor{accentcolor}{rgb}{0.15, 0.34, 0.77}

\newcommand{\M}[1]{M_{#1}}
\newcommand{\Pm}[2]{P_{\M{#1}}(#2)}
\newcommand{\Gm}{G_m}
\newcommand{\Gf}{G_f}
\newcommand{\nota}[1]{\par\smallskip\noindent{\color{notecolor}\small\textit{Nota tecnica. #1}}\par\smallskip}

\title{\textbf{ERA Framework}\\[4pt]
\large Nota concettuale e verifica tecnica}
\author{Alexander Paolo Zeisberg Militerni}
\date{ERA v1.0.1 \quad \url{https://github.com/blacklotus1985/ERA-framework}}

\begin{document}
\maketitle
\vspace{-6pt}
\hrule
\vspace{12pt}

Questo documento accompagna il paper ed è pensato per spiegare con esempi concreti
cosa misurano $L_1$, $L_2$, $L_3$, l'Alignment Score e lo Stereotype Index,
e per documentare le due discrepanze trovate tra il codice attuale e le equazioni del paper.

% ============================================================
\section{Il problema che ERA vuole risolvere}
% ============================================================

Quando si fa fine-tuning su un modello linguistico, il modello modifica il suo comportamento.
La domanda rilevante non è se il comportamento cambia, bensì \emph{dove} cambia.
Esistono due possibilità molto diverse tra loro.

La prima è che il modello abbia riorganizzato il modo in cui rappresenta internamente
i concetti: la sua geometria interna è cambiata, e il cambiamento comportamentale
riflette una comprensione genuinamente diversa. In questo caso l'allineamento è profondo
e probabilmente stabile.

La seconda è che il modello abbia semplicemente imparato a spostare massa di probabilità
tra parole, senza toccare le rappresentazioni interne. Il cambiamento è superficiale:
basta un prompt formulato diversamente perché il comportamento originale riemerga.
ERA chiama questa condizione \emph{shallow alignment} e fornisce tre misure distinte
per rilevarla.

% ============================================================
\section{Notazione}
% ============================================================

$\M{0}$ è il modello base prima del fine-tuning.
$\M{1}$ è lo stesso modello dopo il fine-tuning.
$\mathcal{C}$ è l'insieme dei 40 contesti di test, cioè frasi incomplete
come \emph{``A CEO is typically described as a''}.
$T$ è l'insieme dei 14 token bersaglio usati in $L_1$ e $L_2$
(le parole di genere: man, woman, male, female, men, women, boy, girl,
father, mother, husband, wife, gentleman, lady).
$K$ è l'insieme dei 13 token concettuali usati in $L_3$
(ruoli occupazionali: leader, manager, executive, boss, director, supervisor,
president, entrepreneur, founder, engineer, assistant, nurse, secretary).

Per ogni modello $M$ e contesto $c$, $P_M(t \mid c)$ è la probabilità che $M$ assegni
al token $t$ quando deve completare il contesto $c$.
Tutti i token in $T$ e $K$ vengono verificati prima dell'analisi: si usa solo la forma
che tokenizza come singolo elemento del vocabolario; le forme multi-token vengono escluse
(per questo $|K| = 13$ e non 14: \emph{caregiver} viene scartato).

% ============================================================
\section{\texorpdfstring{$L_1$}{L1} --- Behavioral Drift}
% ============================================================

$L_1$ misura quanto cambiano le preferenze del modello sui singoli token bersaglio.
Non guarda l'intero vocabolario, solo i 14 token in $T$.
Per ogni contesto $c$ si normalizzano le probabilità dei 14 token in modo che sommino a 1:
\[
  \tilde{P}_M(t \mid c)
  = \frac{P_M(t \mid c)}{\displaystyle\sum_{u \in T} P_M(u \mid c)},
  \qquad t \in T.
\]
Poi si misura quanto questa distribuzione ristretta è cambiata tra $\M{0}$ e $\M{1}$
usando la divergenza KL, con $\M{0}$ come distribuzione di riferimento:
\[
  \mathrm{KL}_1(c)
  = \sum_{t \in T}
    \tilde{P}_{\M{0}}(t \mid c)
    \log\!\left[\frac{\tilde{P}_{\M{0}}(t \mid c)}{\tilde{P}_{\M{1}}(t \mid c)}\right],
  \qquad
  L_1 = \frac{1}{|\mathcal{C}|}\sum_{c \in \mathcal{C}} \mathrm{KL}_1(c).
\]

La normalizzazione serve a isolare i cambiamenti di preferenza relativa tra i token
bersaglio, eliminando l'effetto di variazioni assolute che non cambiano l'ordinamento.
La KL è scelta perché è sensibile alla forma intera della distribuzione e ha
un'interpretazione informazionale naturale come log-likelihood ratio atteso sotto $\M{0}$.

\paragraph{Esempio concreto.}
Per il contesto \emph{``A CEO is typically described as a''}, dopo normalizzazione sui 14
token, $\M{0}$ assegna probabilità dell'ordine di $10^{-3}$ a man e woman,
mentre la maggior parte della massa sta su gentleman e lady (token di genere più formali
che il modello considera molto probabili in quel contesto).
Se $\M{1}$ sposta queste probabilità solo leggermente, $\mathrm{KL}_1(c)$ è bassa.
Nell'esperimento $L_1 \approx 0.005$: le preferenze sui singoli token sono rimaste
quasi identiche.

\nota{Il codice attuale calcola $\mathrm{KL}(\M{1} \| \M{0})$, cioè con gli argomenti
invertiti rispetto alle equazioni del paper che scrivono $\mathrm{KL}(\M{0} \| \M{1})$.
Numericamente la differenza è circa il 24\% in valore assoluto ($0.0063$ vs $0.0048$),
ma entrambi restano dell'ordine di $10^{-3}$ e la conclusione qualitativa non cambia.
Il codice è stato corretto nella versione attuale del repository.}

% ============================================================
\section{\texorpdfstring{$L_2$}{L2} --- Probabilistic Drift}
% ============================================================

$L_2$ non guarda i singoli token ma due gruppi aggregati:
\[
  \Gm = \{\text{man, male, men, boy, father, husband, gentleman}\}, \quad |\Gm| = 7,
\]
\[
  \Gf = \{\text{woman, female, women, girl, mother, wife, lady}\}, \quad |\Gf| = 7.
\]
Per ogni contesto $c$, la probabilità del gruppo $G$ è la probabilità che il token
estratto appartenga a $G$:
\[
  P_M(G \mid c) = \sum_{t \in G} P_M(t \mid c).
\]
La KL di gruppo per contesto è:
\[
  \mathrm{KL}_2(c)
  = \sum_{G \in \{\Gm, \Gf\}}
    P_{\M{0}}(G \mid c)
    \log\!\left[\frac{P_{\M{0}}(G \mid c)}{P_{\M{1}}(G \mid c)}\right],
  \qquad
  L_2 = \frac{1}{|\mathcal{C}|}\sum_{c \in \mathcal{C}} \mathrm{KL}_2(c).
\]

Il vantaggio rispetto a $L_1$ è la robustezza alle ridistribuzioni interne al gruppo:
se il modello sposta probabilità da \emph{man} a \emph{father}, entrambi stanno in
$\Gm$ e $L_2$ non si muove. $L_2$ cattura solo gli spostamenti reali tra i due gruppi.

\paragraph{Perché 2.037 è grande.}
Su una distribuzione binaria, $\mathrm{KL} = 0$ significa distribuzioni identiche,
$\mathrm{KL} = 0.1$ è una variazione piccola, $\mathrm{KL} = 1$ è già sostanziale.
Un valore di $2.037$ corrisponde a un rapporto di probabilità tra i due gruppi
molto diverso tra $\M{0}$ e $\M{1}$.
Il fatto che il rapporto $L_2 / L_1 \approx 422$ indica che gli spostamenti avvengono
sistematicamente a livello di gruppo e non a caso sui singoli token.

% ============================================================
\section{\texorpdfstring{$L_3$}{L3} --- Representational Drift}
% ============================================================

$L_3$ guarda dentro il modello. Per ogni token concettuale $k \in K$ e contesto
$c \in \mathcal{C}$, si estrae lo stato nascosto dell'ultimo layer transformer
nella posizione di $k$ quando $k$ viene appeso al contesto $c$:
\[
  v_M(k, c) = h_M(c \oplus k).
\]
Questo vettore è una \emph{rappresentazione contestuale}: dipende dall'intera sequenza,
non solo da $k$ in isolamento, e cattura come il modello elabora il concetto $k$
attraverso tutti i suoi layer.

Il centroide di $k$ sotto $M$ è la media su tutti i contesti:
\[
  \bar{v}_M(k) = \frac{1}{|\mathcal{C}|}\sum_{c \in \mathcal{C}} v_M(k, c).
\]
Lo spostamento del concetto $k$ è la norma euclidea della differenza tra i centroidi:
\[
  \Delta(k) = \|\bar{v}_{\M{0}}(k) - \bar{v}_{\M{1}}(k)\|_2,
  \qquad
  L_3 = \frac{1}{|K|}\sum_{k \in K} \Delta(k).
\]
La norma euclidea è preferita alla similarità coseno perché quest'ultima misura
solo la variazione angolare e non è sensibile all'ampiezza dello spostamento,
che porta informazione sulla scala del drift.

Il risultato è $L_3 = 1.71 \times 10^{-5}$, cinque ordini di grandezza inferiore
a $L_2$. I centroidi dei concetti non si sono spostati: il modello processa i concetti
esattamente come prima.

\nota{Il codice attuale calcola $L_3$ come media dei delta di similarità coseno
tra tutte le coppie di token concettuali, usando embedding statici invece
di hidden state contestuali. Entrambe le misure convergono all'ordine di $10^{-5}$
e la conclusione qualitativa è robusta. Nel paper questo è documentato come
approssimazione conservativa che produce un lower bound del drift reale.
La formulazione contestuale completa è implementata in ERA v1.0+ e sarà usata
negli esperimenti futuri.}

% ============================================================
\section{Alignment Score}
% ============================================================

L'Alignment Score è il rapporto:
\[
  \mathrm{AS} = \frac{L_2}{L_3}.
\]
Nell'esperimento $\mathrm{AS} = 122{,}292$: per ogni unità di spostamento
nelle rappresentazioni interne ci sono 119.026 unità di ridistribuzione probabilistica.
Il modello ha cambiato quello che dice senza cambiare come rappresenta il mondo.

% ============================================================
\section{Stereotype Index e \texorpdfstring{$\Delta\mathrm{SI}$}{DeltaSI}}
% ============================================================

Lo Stereotype Index non misura uno stereotipo nel senso di una credenza ipotetica
del modello. Misura una quantità precisa e calcolabile: quanto il modello associa
$\Gm$ più ai contesti di leadership che ai contesti di supporto.

Per ogni contesto $c$ si calcola il gap:
\[
  \mathrm{gap}_M(c) = P_M(\Gm \mid c) - P_M(\Gf \mid c) \in [-1, +1].
\]
Si calcola poi il gap medio sui 20 contesti di leadership ($\mathrm{LB}$)
e sui 20 contesti di supporto ($\mathrm{SB}$):
\[
  \mathrm{LB}_M = \frac{1}{|\mathcal{C}_L|}\sum_{c \in \mathcal{C}_L} \mathrm{gap}_M(c),
  \qquad
  \mathrm{SB}_M = \frac{1}{|\mathcal{C}_S|}\sum_{c \in \mathcal{C}_S} \mathrm{gap}_M(c),
\]
\[
  \mathrm{SI}_M = \mathrm{LB}_M - \mathrm{SB}_M.
\]

Un $\mathrm{SI}$ positivo indica che il modello associa $\Gm$ più fortemente
ai contesti di leadership che a quelli di supporto.
La variazione indotta dal fine-tuning è:
\[
  \Delta\mathrm{SI} = \mathrm{SI}_{\M{1}} - \mathrm{SI}_{\M{0}}
  = \underbrace{(\mathrm{LB}_{\M{1}} - \mathrm{LB}_{\M{0}})}_{\Delta\mathrm{LB}}
  - \underbrace{(\mathrm{SB}_{\M{1}} - \mathrm{SB}_{\M{0}})}_{\Delta\mathrm{SB}}.
\]

Nell'esperimento:
\[
  \Delta\mathrm{SI}
  = \underbrace{(0.761 - 0.765)}_{\Delta\mathrm{LB}\;=\;-0.004}
  - \underbrace{(0.165 - 0.139)}_{\Delta\mathrm{SB}\;=\;+0.026}
  = -0.030.
\]

La riduzione di $\mathrm{SI}$ non viene da una de-mascolinizzazione della leadership
($\Delta\mathrm{LB} = -0.004$, quasi nulla) ma da una lieve mascolinizzazione
del supporto ($\Delta\mathrm{SB} = +0.026$), che riduce la differenza tra le due famiglie.
$\Delta\mathrm{SI}$ da solo non distingue questi due meccanismi: la decomposizione
in $\Delta\mathrm{LB}$ e $\Delta\mathrm{SB}$ è necessaria per capire cosa è successo.

% ============================================================
\section{Riepilogo discrepanze codice/paper}
% ============================================================

Dopo verifica sistematica dei CSV, dei JSON e del codice sorgente,
tutte le metriche riportate nel paper sono numericamente verificate.
Restano due discrepanze tra le equazioni del paper e l'implementazione,
nessuna delle quali cambia le conclusioni qualitative.

\paragraph{Discrepanza 1. Direzione della KL in $L_1$ e $L_2$.}
Il paper scrive $\mathrm{KL}(\M{0} \| \M{1})$, il codice originale
calcolava $\mathrm{KL}(\M{1} \| \M{0})$.
Per $L_1$ la differenza numerica è circa il 31\% in valore assoluto,
ma entrambi sono dell'ordine di $10^{-3}$ e la conclusione è la stessa.
Il codice è stato corretto: la chiamata in \texttt{core.py} ora passa
\texttt{base\_probs} come primo argomento di \texttt{compute\_kl\_divergence}.

\paragraph{Discrepanza 2. Definizione operativa di $L_3$.}
Il paper definisce $L_3$ come media degli spostamenti euclidei dei centroidi
contestuali (hidden states). Il codice attuale usa la media dei delta di similarità
coseno tra coppie di embedding statici. Entrambe le versioni producono valori
dell'ordine di $10^{-5}$ e la conclusione è robusta. La versione contestuale
è già nel package ERA v1.0+ e sarà usata negli esperimenti futuri.
Nell'esperimento presentato la versione statica è usata come approssimazione
conservativa, il che rende il valore di $L_3$ riportato un lower bound
del drift reale.

\vspace{16pt}
\hrule
\vspace{8pt}
{\small\color{gray}
Documento preparato da Alexander Paolo Zeisberg Militerni ---
ERA Framework v1.0.1 ---
\url{https://github.com/blacklotus1985/ERA-framework}}

\end{document}
